\chapter*{Abstract}
% Ajouter manuellement cette section à la table des matières
\addcontentsline{toc}{chapter}{Abstract}

The \textit{Solvency II} directive requires insurance undertakings to perform a forward-looking assessment of their financial liabilities, necessitating the deployment of stochastic Asset-Liability Management (ALM) models. These projections demand considerable computing power, making the unit-level processing of life insurance contracts operationally prohibitive. To mitigate this algorithmic complexity, reducing the portfolio's dimensionality through an aggregation process into \textit{Model Points} constitutes a suitable solution. However, excessive grouping risks smoothing out financial guarantees, thereby altering the \textit{Best Estimate} (BE) evaluation and the representation of the balance sheet structure under market shocks, potentially distorting the Solvency Capital Requirement (SCR).

The objective of this thesis is to evaluate the influence of different data aggregation methods on the determination of the BE and the SCR. The approach is based on the construction of a reference portfolio of 1,000,000 policies, representative of an average life insurer on the French market.

This portfolio allows for the comparison of several grouping techniques: \textit{Banding}, clustering algorithms (\textit{K-Means} and \textit{HDBSCAN}), as well as a \textit{Cash-Flow Matching} approach. The study compares these methods to identify the one that minimizes the residual error on reference indicators while optimizing computation times.

The analysis tends to demonstrate that dimensionality reduction can decrease processing times by a factor of fifty. In the face of asymmetric interest rate shocks, spatial algorithms—particularly the \textit{K-Means} method—offer the best operational compromise. By preserving the granularity of minimum guaranteed rates and contract seniority, these models accurately capture the financial mechanics of the opportunity cost inherent in lapse risk. Conversely, the approach based on \textit{Cash-Flow Matching} tends to smooth out tax-related information, leading to a mechanical overestimation of the estimated SCR.

\vspace{0.5cm}
\noindent \textbf{Keywords:} life insurance, ALM, Solvency II, aggregation, \textit{Model Points}, \textit{Machine Learning}, \textit{K-Means}, sensitivity analysis, SCR, \textit{Best Estimate}.